\documentclass[11pt]{article}
\usepackage{report}

\title{Nonlinear Band-Edge Thermometry in Periodic Elastic Waveguides:\\
Drift-Resistant Acoustic Temperature Monitoring for Battery Safety}
\author{Kazi Tahsin Mahmood}
\affil{Department of Mechanical Engineering, Wayne State University}
\date{}

\begin{document}

\maketitle

\section{Intellectual Merit}

\paragraph{Challenge.}
Lithium-ion battery thermal runaway is a persistent safety risk in transportation, grid storage, and defense. Interior hot spots can precede external warning signs by minutes to hours, yet most commercial battery management systems rely on surface temperature measurements; direct internal measurement generally requires embedded or penetrative hardware. Ultrasonic time-of-flight (TOF) is promising for non-invasive thermometry~\cite{Cheng2024,Williams2024}, but TOF features must be decoupled from state of charge, coupling condition, gas evolution, and mechanical boundary changes. This project addresses that gap by combining nonlinear periodic-structure amplification with phase-sensitive acoustic readout.

\paragraph{Why This is New and Why Now.}
Linear ultrasonic battery thermometry exists~\cite{Cheng2024,Williams2024}; however, nonlinear band-edge amplification of second-harmonic generation remains largely unexplored as a battery temperature discriminator. The innovation is converting a small temperature-induced stiffness shift into a nonlinear amplitude-ratio signal through engineered dispersion. The thermometric waveguide operates near a band edge, where slow group velocity localizes energy and the second-harmonic response is expected to exhibit inverse detuning dependence. Because the measured response contains amplitude and phase, the band-edge field can also be treated as a complex acoustic state whose temperature-induced geometric phase change provides the topological-acoustic diagnostic~\cite{Zhang2024SensorsGPC}. In this framing, the nonlinear band edge amplifies the thermal perturbation, while the phase-geometric metric tests whether that perturbation remains visible in the global acoustic state rather than only in a local amplitude. Thus, the topological-acoustic element is the geometric-phase description of how the complex band-edge acoustic state evolves under thermal perturbation, complementing the nonlinear harmonic-ratio readout. In the reduced model, harmonic balance gives $A_{2f}/A_f \sim \delta_\mathrm{NL}/(\Delta^2+\zeta^2)$, where $\Delta = f_\mathrm{drive} - f_\mathrm{edge}(T)$, $\delta_\mathrm{NL}$ is the effective nonlinear coefficient, and $\zeta$ represents damping/linewidth. This project will test whether shifts in $f_\mathrm{edge}(T)$ produce harmonic-ratio and phase-geometric temperature metrics beyond conventional wave-speed readouts. The timing is right: nonlinear periodic-structure theory supports band-edge-sensitive design~\cite{Lazarov2007}, CNC machining enables reproducible asymmetric stubs, and battery deployment creates urgent demand for interior temperature monitoring beyond surface-only BMS measurements.

\paragraph{PI Background.}
The PI's research expertise spans nonlinear vibration, acoustics, periodic elastic structures, and phase-sensitive state reconstruction. Prior work demonstrates phase-sensitive lock-in detection of nonlinear harmonic states in piezo-driven acoustic waveguides~\cite{Mahmood2022} and a Bloch-sphere dispersion formalism for one-dimensional elastic lattices covering band-edge analysis and symmetry-governed mode hybridization~\cite{Mahmood2026a}. These contributions provide direct preparation through harmonic balance, Floquet-Bloch dispersion, multiple-scales perturbation, and complex acoustic-state analysis.


\section{Research Plan}

\paragraph{Tasks.}
\begin{itemize}[noitemsep]
  \item \textbf{Task~1: Analytical model (months 1-4):} \emph{Problem:} A predictive model is needed that links the unit-cell Duffing coefficient, detuning $\Delta$, damping/linewidth $\zeta$, temperature sensitivity of $A_{2f}/A_f$, and the geometric-phase-change index. \emph{Solution:} Harmonic balance at $f_0$ and $2f_0$ applied to a Duffing diatomic chain will derive a leading-order expression expected to scale as $A_{2f}/A_f \sim \delta_\mathrm{NL}/(\Delta^2+\zeta^2)$ under weakly nonlinear, near-band-edge assumptions. Multiple-scales analysis will bound the valid drive-amplitude range and predict a steep but bounded temperature derivative, with $\Delta^{-3}$ recovered only as a low-damping asymptote. The deliverable is a design map for the useful detuning window, required nonlinear coefficient, damping tolerance, signal-to-drift margin, and baseline-to-perturbed geometric phase change. The geometric-phase-change index will be computed from the normalized complex response field at selected receiver locations or modal components, identifying regimes where the phase-geometric metric and harmonic ratio remain distinguishable from ordinary drive-level or coupling variations.

  \item \textbf{Task~2: FE simulation (months 3-8):} \emph{Problem:} The 1D chain model neglects plate boundaries, two-dimensional spreading, and higher-order modes that can contaminate the $2f$ channel. \emph{Solution:} A 30-40\,cm 6061-T6 aluminum plate waveguide with 8-12 CNC-asymmetric unit cells is modeled with temperature-dependent elastic constants. Stub-height, width, and asymmetry sweeps map $\partial(A_{2f}/A_f)/\partial T$; $3f$/$4f$ monitoring identifies a clean $2f$ regime not dominated by broadband nonlinear distortion. Dispersion and mode-shape extraction verify that the selected drive excites the intended band-edge mode rather than a plate resonance artifact. A symmetric-stub FE control, fabrication-tolerance check, and 10\% agreement target with Task~1 will benchmark the reduced model and fix the final geometry before machining.

  \item \textbf{Task~3: Fabrication and experiment (months 6-12):} \emph{Problem:} Source harmonics at $2f_0$ and coupling-layer drift can mimic the structural temperature signal. \emph{Solution:} A transmit bandpass filter suppresses source harmonics; $A_{2f}/A_f$ is extracted by coherent FFT averaging over 20 tone-bursts per step as the specimen is stepped through 20--100$^\circ$C in 5$^\circ$C increments, with drive voltage modulated $\pm$20\% to emulate coupling variation. A symmetric-stub control, off-band drive point, repeated coupling perturbation, and simultaneous linear TOF measurement separate bond, amplifier, and wave-speed effects from the engineered nonlinear response. Sensitivity, repeatability, drift-induced temperature error, and geometric-phase-change temperature dependence are reported with 95\% confidence intervals from ten repeated cycles, providing a direct comparison among nonlinear band-edge thermometry, phase-geometric diagnostics, and conventional TOF. The geometric-phase trend will be checked for monotonicity and repeatability across temperature cycles to determine whether it provides an independent phase-based indicator of thermal perturbation.
\end{itemize}

\paragraph{Milestones, risk, and timeline.}
\textbf{M1 (month 4):} Reduced-order sensitivity model and analytical design map completed. \textbf{M2 (month 8):} FE sensitivity map completed; results presented at a conference. \textbf{M3 (month 12):} Plate specimen fabricated and tested; sensitivity, drift-rejection, and repeatability data collected; journal manuscript submitted. The primary risk is that stub nonlinearity is too small to produce a measurable second-harmonic ratio or resolved phase-geometric change; mitigation includes band-edge operation, symmetric-stub control, geometry optimization before machining, and input harmonic filtering. \textbf{Timeline:} Sept--Dec~2026 analytical modeling; Jan--Apr~2027 FE simulation and fabrication; May--Aug~2027 experiments, data analysis, and manuscript writing.


\section{Broader Impact}

\paragraph{Technological.}
This project aims to establish acoustic thermometry in which temperature sensitivity is amplified by nonlinear band-edge dynamics and interpreted through harmonic-ratio and phase-geometric acoustic-state metrics, rather than direct wave-speed measurement alone. If validated, the framework could generalize to strain, state-of-charge, electrolyte degradation, or early hot-spot detection by redesigning the unit-cell geometry and selecting the acoustic-state metric most responsive to the target perturbation. Successful demonstration would position the work for follow-on support in nonlinear dynamics, energy-storage safety, and acoustic diagnostics, with Year-2 validation against commercial Li-ion cells under charge-discharge cycling.

\paragraph{Workforce and education.}
This project will develop the PI's expertise in nonlinear vibration, periodic elastic structures, topological acoustic diagnostics, and energy-storage safety. Undergraduate research opportunities will be targeted at NewFoS partner Spelman College, and models, scripts, drawings, and calibration datasets will be released through the NewFoS community data portal.

{%
\renewcommand*{\bibfont}{\scriptsize}
\setlength{\bibitemsep}{1pt}
\setlength{\bibnamesep}{0pt}
\printbibliography
}

\end{document}
